\documentclass[a4paper,12pt]{article}

% Pacchetti utili
\usepackage[italian]{babel}
\usepackage[T1]{fontenc}
\usepackage[utf8]{inputenc}
\usepackage{lmodern}
\usepackage{geometry}
\usepackage{setspace}
\usepackage{graphicx}
\usepackage{titling}
\usepackage{hyperref}
\usepackage{booktabs}
\usepackage{amsmath, amssymb}
\usepackage[most]{tcolorbox}
\usepackage{tikz}

\geometry{margin=2.5cm}
\onehalfspacing

\begin{document}
	
	\begin{titlepage}
		\centering
		
		% Intestazione Università
		{\scshape\LARGE Università di Pisa \par}
		\vspace{0.5cm}
		
		% Titolo corso
		{\Huge\bfseries Calcolo Numerico - Esercizi \par}
		\vspace{0.5cm}
		{\Large Anno Accademico 2025--2026 \par}
		
		\vfill
		
		% Autore
		{\Large\textbf{Pietro Bacigalupi} \par}
		\vspace{0.5cm}
		
		\vfill
		
		% Data
		{\large \today\par}
		
	\end{titlepage}
	
	\tableofcontents
	\newpage
	
	\section{Calcolo degli Errori}
	\paragraph{TODO...}
	
	\section{Matrici}
	\subsection{Esercizio 1 (Norme) }
		Data la matrice
		\begin{align*}
			A = \begin{bmatrix}
				2 & -1 \\
				4 & -3
			\end{bmatrix}
		\end{align*}
		Calcola:
		\begin{enumerate}
			\item $\|A\|_1$
			\item $\|A\|_{\infty}$
			\item $\|A\|_F$ (Norma di Frobenius)
		\end{enumerate}
		\paragraph{Risoluzione}
		\begin{enumerate}
			\item $\|A\|_1 = \max\limits_{j=1,...,n} \sum_{i=1}^{n} |a_{ij}|$ (Massima somma dei moduli degli elementi di ogni colonna).
			
			\begin{itemize}
				\item Colonna 1: $j = 1: |2| + |4| = 6$
				\item Colonna 2: $j = 2: |-1| + |-3| = 4$
			\end{itemize}
			Prendo il massimo tra $6$ e $4$: $\|A\|_1 = 6$
			\item $\|A\|_{\infty} = \max\limits_{i=1,...,n} \sum_{j=1}^{n} |a_{ij}|$ (Massima somma dei moduli degli elementi di ogni riga).
			
			\begin{itemize}
				\item Riga 1: $i = 1: |2| + |-1| = 3$
				\item Riga 2: $i = 2: |4| + |-3| = 7$
			\end{itemize}
			Prendo il massimo tra $3$ e $7$: $\|A\|_{\infty} = 7$
			
			\item $\|A\|_F = \sqrt{\sum_{i,j = 1}^{n} |a_{ij}|^2}$ (Radice della somma dei quadrati di tutti gli elementi della matrice.)
			
			$|2|^2 + |-1|^2 + |4|^2 + |-3|^2 = 30$.
			
			 quindi estraendo la radice otteniamo: $\|A\|_F = \sqrt{30}$
		\end{enumerate}
		\subsection{Esercizio 2 {Gershgorin}}
		Data la matrice:
		\begin{align*}
			\begin{bmatrix}
				5 & 1 & -1 \\
				2 & -4 & 0 \\
				0 & 1 & 2 
			\end{bmatrix}
		\end{align*}
		\begin{enumerate}
			\item Determina il centro $c_{i}$ e il raggio $\rho_i$ dei tre cerchi di Gershgorin associati alle righe della matrice B
			\item Verificando le disuguaglianze, stabilisci se la matrice gode della proprietà di predominanza diagonale forte.
			\item Sulla base della risposta precedente, possiamo garantire al 100\% che la matrice sia invertibile?
		\end{enumerate}
		\paragraph{Risoluzione}
		
		\begin{enumerate}
			\item Calcoliamo il centro di ogni cerchio: $\forall(i=1,...,n), c_i =$ Elemento della diagonale principale di quella riga.
			\begin{itemize}
				\item $[5, 1, -1]$ $c_1 = 5$
				\item $[2, -4, 0]$ $c_2 = -4$
				\item $[0, 1, 2]$ $c_3 = 2$
			\end{itemize}
			Il raggio $\rho_i = \sum_{j\neq i}^{n} |a_{ij}|$ è dato dalla somma dei moduli di ogni elemento della riga diversi da $c_i$.
			\begin{itemize}
				\item $[5, 1, -1]$ $\rho_1 = |1| + |-1| = 2$
				\item $[2, -4, 0]$ $\rho_2 = |2| + |0| = 2$
				\item $[0, 1, 2]$ $\rho_3 = |0| + |1| = 1$
			\end{itemize} 
			\item Una matrice gode della proprietà di predominanza diagonale forte se, per ogni riga, il modulo dell'elemento sulla diagonale principale è strettamente maggiore della somma dei moduli degli altri elementi della riga ($\rho_i$).
			
			Verifichiamo la condizione $|c_i| > \rho_i$ per ogni riga:
			\begin{itemize}
				\item \textbf{Riga 1}: $|5| > 2$ (Vero)
				\item \textbf{Riga 2}: $|-4| > 2$ (Vero)
				\item \textbf{Riga 3}: $|2| > 1$ (Vero)
			\end{itemize}
			Poiché la disuguaglianza stretta è verificata per tutte e tre le righe, possiamo affermare che \textbf{la matrice B è a predominanza diagonale forte}
			\item Sì, possiamo garantirlo al 100\%. Se una matrice è a predominanza diagonale forte, allora è sicuramente invertibile. Questo perché i raggi dei cerchi di Gershgorin sono rigorosamente minori della distanza dei rispettivi centri dall'origine del piano complesso (0). Di conseguenza, lo zero non cade in nessun cerchio di Gershgorin e, per il $1^o$ teorema, non potendo essere un autovalore, il determinante della matrice è diverso da zero $(det(B)\neq0)$.
		\end{enumerate}
		
	
\end{document}
